\subsection{分歧指数. 剩余类次数}

设 K 是域，L 是 K 的扩张，A' 是 L 的赋值环。如在 § 1 第 4 节所见，环 A = K ∩ A' 是 K 的赋值环，且

\[
\mathfrak {m} (\mathrm{A}) = \mathfrak {m} (\mathrm{A} ^ {\prime}) \cap \mathrm{K}.
\]

若 \(v'\) 是与 \(A'\) 相伴的赋值，则 \(v'\) 在 K 上的限制 v 是与 A 相伴的 K 上的赋值；v 的阶群 \(\Gamma_{v}\) 是 \(\Gamma_{v'}\) （\(v'\) 的阶群）的子群。

定义 1. 指标 \((\Gamma_v: \Gamma_{v'})\) 称为 \(v'\) 在 \(v\) 上（或在 K 上）的分歧指数，记作 \(e(v'/v)\) （或 \(e(A'/A)\) ，有时也记作 \(e(L/K)\) ）。

这个指标是自然数或 \(+\infty\) 。若 \(v_0'\) 是与 \(v'\) 等价的赋值，\(e(v'/v)\) 也称为 \(v_0'\) 在 \(v\) 上的分歧指数。若 \(e(v'/v) = 1\) ，则称 \(v'\) 在 \(v\) 上非分歧。

另一方面，剩余域 \(\kappa(A)\) （\(v\) 的）与剩余域 \(\kappa(A')\) （\(v'\) 的）的一个子域视为同一。

定义 2. 次数 \([\kappa (\mathbf{A}^{\prime}):\kappa (\mathbf{A})]\) 称为 \(v^{\prime}\) 在 \(v\) 上（或在 K 上）的剩余类次数，记作 \(f(v^{\prime} / v)\) （或 \(f(\mathbf{A}^{\prime} / \mathbf{A})\) ，有时也记作 \(f(\mathbf{L} / \mathbf{K})\) ）。

这个次数是自然数或 +∞。

引理 1. 设 K 、 K' 、 K'' 是三个域，满足 K ⊂ K' ⊂ K'' ，v'' 是 K'' 上的赋值，v 与 v' 是它在 K 与 K' 上的限制。那么有下列关系：

\[
e (v ^ {\prime \prime} / v) = e (v ^ {\prime \prime} / v ^ {\prime}) e (v ^ {\prime} / v), \quad f (v ^ {\prime \prime} / v) = f (v ^ {\prime \prime} / v ^ {\prime}) f (v ^ {\prime} / v).\tag{1}
\]

这是显然的。

引理 2. 设 K 是域，L 是 K 的次数为 n 的有限扩张，v' 是 L 上的赋值，v 是它在 K 上的限制。那么不等式

\[
e (v ^ {\prime} / v) f (v ^ {\prime} / v) \leqslant n\tag{2}
\]

成立；特别地 \(e(v' / v)\) 与 \(f(v' / v)\) 有限。

取自然数 \(r\) 与 \(s\) ，分别不大于 \(e(v'/v)\) 与 \(f(v'/v)\) 。只需证明 \(rs \leqslant n\) 。鉴于 \(r\) 的定义，存在元素 \(x_i\) 属于 \(L\) ( \(1 \leqslant i \leqslant r\) ) 使得 \(v'(x_i) \not\equiv v'(x_j) \pmod{\Gamma_v}\) 对 \(i \neq j\) 成立。鉴于 \(s\) 的定义，存在元素 \(y_k\) ( \(1 \leqslant k \leqslant s\) ) 属于环 \(A'\) （\(v'\) 的环），它们的典范像 \(\overline{y}_k\) 在 \(K(A')\) 中在 \(K(A)\) 上线性无关；显然 \(v'(y_k) = 0\) 对一切 \(k\) 成立。我们将证明这 \(rs\) 个元素 \(x_i y_k\) 在 \(K\) 上线性无关，这当然就建立了不等式 \(rs \leqslant n\) 。

于是设存在一个形如下式的非平凡线性关系

\[
\sum_ {i, k} a _ {i k} x _ {i} y _ {k} = 0 \quad (a _ {i k} \in \mathbf {K}).\tag{3}
\]

选取指标 \(j, m\) 使得

\[
v ^ {\prime} (a _ {j m} x _ {j} y _ {m}) \leqslant v ^ {\prime} (a _ {i k} x _ {i} y _ {k})
\]

对每个有序对 \((i,k)\) 成立；那么 \(a_{jm} \neq 0\) 。若 \(i \neq j\) ，则 \(v'(a_{ik}x_iy_k) = v'(a_{jm}x_jy_m)\) 不可能成立，因为那将蕴含

\[
v ^ {\prime} (x _ {i}) - v ^ {\prime} (x _ {j}) = v ^ {\prime} (a _ {j m}) - v ^ {\prime} (a _ {i k}) \in \Gamma_ {v},
\]

这与诸 \(x_{i}\) 的选取矛盾。用 \((a_{jm}x_j)^{-1}\) 乘 (3) ，得到一个形如下式的关系

\[
\sum_ {k} b _ {k} y _ {k} + z = 0, \quad \text { where } \quad b _ {k} = \frac {a _ {j k} x _ {j}}{a _ {j m} x _ {j}} \in \mathrm{A} ^ {\prime}, \qquad z \in \mathrm{A} ^ {\prime}
\]

且 \(v'(b_k) \geqslant 0, v'(z) > 0\) 。于是在 \(\kappa(A')\) 中得到形如 \(\sum_{k} \bar{b}_k \bar{y}_k = 0\) 的关系。由于 \(b_m = 1\) ，这与对 \(y_k\) 所作的假设矛盾。

命题 1. 设 K 是域，L 是 K 的代数扩张，\(v'\) 是 L 上的赋值，\(v\) 是它在 K 上的限制，A 与 A' 是 \(v\) 与 \(v'\) 的环。那么 \(\Gamma_{v'} / \Gamma_v\) 是扭群，且 \(\kappa(A')\) 是 \(\kappa(A)\) 的代数扩张。

设 \((\mathbf{L}_{\alpha})\) 是 L 的有限子扩张的族；记 \(\Gamma_{\alpha} = v'(\mathbf{L}_{\alpha}^{*})\) 。群 \(\Gamma_{v'}\) 是由诸 \(\Gamma_{\alpha}\) 组成的右向定向族的并；由于诸群 \(\Gamma_{\alpha}/\Gamma_{v}\) 有限（引理 2），\(\Gamma_{v'}/\Gamma_{v}\) 是扭群。为证明 \(\kappa(\mathbf{A}')\) 是 \(\kappa(\mathbf{A})\) 的代数扩张，论证类似。

推论 1. \(v'\) 的高等于 \(v\) 的高。

这由命题 1 与下面的引理得出：

引理 3. 设 G' 是全序群，G 是 G' 的子群，S' （分别 S）是 G' （分别 G）的孤立子群集。映射 \(H' \mapsto H' \cap G\) 把 S' 映到 S 上。若 G'/G 是扭群，这个映射是双射。

显然 \(H' \in S'\) 蕴含 \(H' \cap G \in S\) 。现在设 \(H \in S\) ；用 \(H'\) 表示由这样的 \(x' \in G'\) 组成的集合：存在 \(h \in H\) 满足 \(-h \leqslant x' \leqslant h\) ；立即可以验证 \(H'\) 是 \(G'\) 的孤立子群；于是由于 H 孤立，有 \(H' \cap G = H\) ；因而映射 \(H' \mapsto H' \cap G\) 是满射。最后设 \(G'/G\) 是扭群；设 \(H_1'\) 与 \(H_2'\) 是 \(G'\) 的两个孤立子群，满足 \(H_1' \cap G = H_2' \cap G\) ；那么例如 \(H_1' \supset H_2'\) （参见~§4 第 4 节）；于是 \(H_1'/H_2'\) 是全序群，且同构于 \(H_1'/(H_1' \cap G)\) 的一个商群，而后者本身又与 \(G'/G\) 的一个子群视为同一；因而 \(H_1'/H_2'\) 是扭群，从而化为 0。

推论 2. \(v'\) 是非真的（分别高为 1），当且仅当 \(v\) 是非真的（分别高为 1）。

推论 3. 设 L 是 K 的有限扩张。\(v'\) 是离散的，当且仅当 \(v\) 是离散的。

若 \(v'\) 是离散的，则 \(\Gamma_v\) 同构于 \(\mathbf{Z}\) 的非零子群（推论 2），因而同构于 \(\mathbf{Z}\) 。反之，若 \(v\) 是离散的，则 \(\Gamma_v\) 同构于 \(\mathbf{Z}\) ，且 \(\Gamma_{v'} / \Gamma_v\) 是有限群（引理 2）；因而 \(\Gamma_{v'}\) 是秩 1 的无扭有限生成交换群；从而它同构于 \(\mathbf{Z}\) 。

